\section{Well-Ordering Principle}

\begin{frame}
\centering
\Huge\textbf{Well-Ordering Principle}
\end{frame}

% --------------------------------------------------
% Question 1
% --------------------------------------------------

\begin{frame}{Question 1}
A software system assigns positive integer IDs to tasks.

Suppose a collection of pending tasks is represented by a non-empty subset

$$
S\subseteq\mathbb{Z}^+.
$$

\textbf{Question:} What property guarantees that there is a task with the smallest ID?
\end{frame}

\begin{frame}{Answer 1}
The required property is the \textbf{Well-Ordering Principle}.

It states that every non-empty subset of the positive integers contains a least element.

Since

$$
S\subseteq\mathbb{Z}^+
$$

and $S$ is non-empty,

$$
\boxed{\exists m\in S\text{ such that }m\leq x
\text{ for every }x\in S.}
$$

Therefore, the pending tasks have a smallest ID.
\end{frame}

% --------------------------------------------------
% Question 2
% --------------------------------------------------

\begin{frame}{Question 2}
A distributed system records retry attempts using positive integer numbers.

Let $S$ be the set of retry numbers for which a particular event succeeds.

Suppose at least one successful retry exists.

\textbf{Question:} Explain why there must be a first successful retry.
\end{frame}

\begin{frame}{Answer 2}
Since at least one successful retry exists,

$$
S\neq\varnothing.
$$

Also,

$$
S\subseteq\mathbb{Z}^+.
$$

By the Well-Ordering Principle, every non-empty subset of positive integers has a least element.

Therefore, there exists
$$
m\in S
$$

such that
$$
m\leq x
\qquad\forall x\in S.
$$
Hence, $m$ is the first successful retry.
\end{frame}

% --------------------------------------------------
% Question 3
% --------------------------------------------------

\begin{frame}{Question 3}
Consider the set

$$
S=\{n\in\mathbb{Z}^+\mid n^2>20\}.
$$

A program searches for the smallest positive integer satisfying the condition.

\textbf{Question:} Use the Well-Ordering Principle to justify why a smallest valid value exists.
\end{frame}

\begin{frame}{Answer 3}
First, observe that $S$ is non-empty because

$$
5^2=25>20.
$$

Therefore,

$$
5\in S.
$$

Also,

$$
S\subseteq\mathbb{Z}^+.
$$

By the Well-Ordering Principle, $S$ must contain a least element.
\end{frame}

\begin{frame}{Answer 3}
We can identify the least element directly.

For positive integers less than $5$,

$$
1^2=1,\quad
2^2=4,\quad
3^2=9,\quad
4^2=16.
$$

None is greater than $20$.

But

$$
5^2=25>20.
$$

Therefore,

$$
\boxed{\min(S)=5}.
$$

\end{frame}

% --------------------------------------------------
% Question 4
% --------------------------------------------------

\begin{frame}{Question 4}
A programmer claims:

\textit{Every non-empty subset of the integers has a smallest element.}

\textbf{Question:} Is this statement true? If not, provide a counterexample and explain why the Well-Ordering Principle does not apply.
\end{frame}

\begin{frame}{Answer 4}
The statement is false for arbitrary subsets of the integers.

Consider

$$
S=\{\ldots,-3,-2,-1\}.
$$

This set is non-empty, but it has no smallest element.

For any

$$
n\in S,
$$

the integer
$$
n-1
$$

also belongs to $S$ and satisfies
$$
n-1<n.
$$

\end{frame}

\begin{frame}{Answer 4}
The Well-Ordering Principle applies specifically to non-empty subsets of the positive integers.

It does not state that every non-empty subset of $\mathbb{Z}$ has a least element.

Therefore,

$$
\boxed{\text{The programmer's claim is false.}}
$$

\end{frame}

% --------------------------------------------------
% Question 5
% --------------------------------------------------

\begin{frame}{Question 5}
A search algorithm examines positive integer candidate values and succeeds for at least one candidate.

Let $S$ be the set of successful candidate values.

The algorithm always chooses the smallest successful candidate.

\textbf{Question:} Explain mathematically why this strategy is well-defined.
\end{frame}

\begin{frame}{Answer 5}
Since the algorithm succeeds for at least one candidate,

$$
S\neq\varnothing.
$$

Every candidate is a positive integer, so

$$
S\subseteq\mathbb{Z}^+.
$$

By the Well-Ordering Principle, $S$ contains a least element.

Therefore,

$$
\exists m\in S
$$

such that

$$
m\leq x
\qquad\forall x\in S.
$$

\end{frame}

\begin{frame}{Answer 5}
Hence, the instruction

\textit{choose the smallest successful candidate}

is mathematically well-defined.

The algorithm can select

$$
\boxed{m=\min(S)}.
$$

The existence of this minimum follows from the Well-Ordering Principle.
\end{frame}
